\documentclass[12pt]{article}
\usepackage[a4paper,margin=1in]{geometry}
\usepackage[utf8]{inputenc}
\usepackage[T1]{fontenc}
\usepackage{amsmath,amssymb}
\usepackage{textcomp}
\usepackage{parskip}
\begin{document}
\begin{center}
{\Large \textbf{2003 STEP III}}
\end{center}

\subsection*{Page 1}
\begin{flushleft}
STEP III, 2003 2\par
\bigskip\textbf{Section A: Pure Mathematics}\par
\bigskip\textbf{Question 1.} Given that x + a > 0 and x + b > 0 , and that b > a , show that\par
\[d x + a \sqrt\{b\} a\]
\[arcsin = -\]
\[dx \lambda x + b \mu (x + b) \sqrt\{a\} + b + 2x\]
d x + b\par
and find arcosh .\par
dx \lambda x + a \mu\par
\[Hence, or otherwise, integrate, for x > 1 ,\]
-\par
\textbf{(i)} dx ,\par
\[Z (x + 1) \sqrt\{x\} + 3\]
\textbf{(ii)} dx .\par
\[Z (x + 3) \sqrt\{x\} + 1\]
d 1 d 1\par
\[[You may use the results arcsin x = and arcosh x = . ]\]
\[dx \sqrt\{1\} x^2 dx \sqrt\{x^2\} 1\]
- -\par
1 3 (2r 1)\par
\bigskip\textbf{Question 2.} Show that 2r Cr = \times \times \cdots \times - 2 r , for r > 1 .\par
\[r! \times\]
- 1\par
\textbf{(i)} Give the first four terms of the binomial series for (1 p) 2 .\par
-\par
By choosing a suitable value for p in this series, or otherwise, show that\par
\infty 2r\par
Cr\par
\[= \sqrt\{2\} .\]
8r\par
\[Xr=0\]
\textbf{(ii)} Show that\par
\infty\par
(2r + 1)\par
2r\par
Cr 3\par
\[= \sqrt\{5\} .\]
5r\par
\[Xr=0\]
n\par
\[[Note: n Cr is an alternative notation for for r > 1 , and 0 C^0 = 1 .]\]
\lambdar\mu\par
\end{flushleft}
\subsection*{Page 2}
\begin{flushleft}
STEP III, 2003 3\par
\[3 If m is a positive integer, show that (1 + x) m + (1 x) m = 0 for any real x .\]
- 6\par
The function f is defined by\par
m m\par
(1 + x) (1 x)\par
\[f (x) = - - .\]
(1 + x)m + (1 x)m\par
-\par
Find and simplify an expression for f (x).\par
\[In the case m = 5 , sketch the curves y = f (x) and y = .\]
f (x)\par
\bigskip\textbf{Question 4.} A curve is defined parametrically by\par
\[x = t^2 , y = t(1 + t^2 ) .\]
The tangent at the point with parameter t, where t = 0 , meets the curve again at the point\par
\[with parameter T , where T = t . Show that\]
\[1 t^2\]
\[T = - and 3t2 = 1 .\]
2t 6\par
Given a point P 0 on the curve, with parameter t 0 , a sequence of points P 0 , P 1 , P 2 , . . . on the\par
curve is constructed such that the tangent at Pi meets the curve again at Pi+1. If t 0 = tan 7 1 \pi 8 ,\par
show that P 3 = P 0 but P 1 = P 0 . Find a second value of t 0 , with t 0 > 0 , for which P 3 = P 0\par
\[but P 1 = P 0 .\]
\bigskip\textbf{Question 5.} Find the coordinates of the turning point on the curve y = x2 2bx + c . Sketch the curve in\par
the case that the equation x2 2bx + c = 0 has two distinct - real roots. Use your sketch to\par
determine necessary and suffici - ent conditions on b and c for the equation x2 2bx + c = 0 to\par
-\par
have two distinct real roots. Determine necessary and sufficient conditions on b and c for this\par
equation to have two distinct positive roots.\par
Find the coordinates of the turning points on the curve y = x3 3b2x+c (with b > 0) and hence\par
determine necessary and sufficient conditions on b and c for th - e equation x3 3b2x + c = 0 to\par
-\par
have three distinct real roots. Determine necessary and sufficient conditions on a , b and c for\par
\[the equation (x a) 3 3b2 (x a) + c = 0 to have three distinct positive roots.\]
- - -\par
\[Show that the equation 2x3 9x2 + 7x 1 = 0 has three distinct positive roots.\]
- -\par
\end{flushleft}
\subsection*{Page 3}
\begin{flushleft}
STEP III, 2003 4\par
\bigskip\textbf{Question 6.} Show that\par
\[2 sin 1 θ cos rθ = sin r + 1 θ sin r 1 θ .\]
2 2 - - 2\par
Hence, or otherwise, find all solutions of the equation\par
\[cos aθ + cos(a + 1)θ + + cos(b 2)θ + cos(b 1)θ = 0 ,\]
\cdots - -\par
\[where a and b are positive integers with a < b 1 .\]
-\par
\bigskip\textbf{Question 7.} In the x-y plane, the point A has coordinates (a , 0) and the point B has coordinates (0 , b) ,\par
where a and b are positive. The point P , which is distinct from A and B, has coordinates (s, t) .\par
X and Y are the feet of the perpendiculars from P to the x-axis and y-axis respectively, and\par
N is the foot of the perpendicular from P to the line AB . Show that the coordinates (x , y)\par
of N are given by\par
\[ab2 a(bt as) a2b + b(bt as)\]
\[x = - - , y = - .\]
\[a^2 + b^2 a^2 + b^2\]
t b t\par
\[Show that, if - = 1 , then N lies on the line XY .\]
\lambda s \mu \lambda s a \mu -\par
-\par
Give a geometrical interpretation of this result.\par
\bigskip\textbf{Question 8.} (i) Show that the gradient at a point (x , y) on the curve\par
\[(y + 2x) (y 4x) = c ,\]
-\par
where c is a constant, is given by\par
dy 16x y\par
\[= - .\]
dx 2y 5x\par
-\par
n\par
\textbf{(ii)} By considering the derivative with respect to x of (y + ax) (y + bx) , or otherwise, find\par
the general solution of the differential equation\par
dy 10x 4y\par
\[= - .\]
dx 3x y\par
-\par
\end{flushleft}
\subsection*{Page 4}
\begin{flushleft}
STEP III, 2003 5\par
\bigskip\textbf{Section B: Mechanics}\par
\bigskip\textbf{Question 9.} A particle P of mass m is constrained to move on a vertical circle of smooth wire with centre O\par
\[and of radius a. L is the lowest point of the circle and H the highest and ∠LOP = θ . The\]
particle is attached to H by an elastic string of natural length a and modulus of elasticity αmg ,\par
\[where α > 1 . Show that, if α > 2 , there is an equilibrium position with 0 < θ < \pi .\]
\[\pi\]
\[Given that α = 2 + \sqrt\{2\} , and that θ = + φ , show that\]
\[g(\sqrt\{2\} + 1)\]
φ¨ φ\par
\approx - 2a\par
when φ is small.\par
For this value of α, explain briefly what happens to the particle if it is given a small displace-\par
\[\pi\]
\[ment when θ = .\]
\bigskip\textbf{Question 10.} A particle moves along the x-axis in such a way that its acceleration is kxx˙ where k is a\par
\[positive constant. When t = 0, x = d (where d > 0) and x˙ = U .\]
\textbf{(i)} Find x as a function of t in the case U = kd2 and show that x tends to infinity as\par
\[\pi\]
t tends to .\par
2dk\par
\textbf{(ii)} If U < 0, find x as a function of t and show that it tends to a limit, which you should\par
state in terms of d and U , as t tends to infinity.\par
11 Point B is a distance d due south of point A on a horizontal plane. Particle P is at rest at\par
B at t = 0, when it begins to move with constant acceleration a in a straight line with fixed\par
bearing β . Particle Q is projected from point A at t = 0 and moves in a straight line with\par
constant speed v . Show that if the direction of projection of Q can be chosen so that Q strikes\par
P , then\par
\[v^2 > ad (1 cos β) .\]
-\par
Show further that if v2 > ad(1 cos β) then the direction of projection of Q can be chosen so\par
-\par
that Q strikes P before P has moved a distance d .\par
\end{flushleft}
\subsection*{Page 5}
\begin{flushleft}
STEP III, 2003 6\par
\bigskip\textbf{Section C: Probability and Statistics}\par
12 Brief interruptions to my work occur on average every ten minutes and the number of inter-\par
ruptions in any given time period has a Poisson distribution. Given that an interruption has\par
just occurred, find the probability that I will have less than t minutes to work before the next\par
interruption. If the random variable T is the time I have to work before the next interruption,\par
find the probability density function of T .\par
I need an uninterrupted half hour to finish an important paper. Show that the expected\par
number of interruptions before my first uninterrupted period of half an hour or more is e3 1.\par
-\par
Find also the expected length of time between interruptions that are less than half an hour\par
apart. Hence write down the expected wait before my first uninterrupted period of half an\par
hour or more.\par
\bigskip\textbf{Question 13.} In a rabbit warren, underground chambers A, B, C and D are at the vertices of a square, and\par
burrows join A to B, B to C, C to D and D to A. Each of the chambers also has a tunnel\par
to the surface. A rabbit finding itself in any chamber runs along one of the two burrows to a\par
neighbouring chamber, or leaves the burrow through the tunnel to the surface. Each of these\par
three possibilities is equally likely.\par
Let pA , pB , pC and pD be the probabilities of a rabbit leaving the burrow through the tunnel\par
from chamber A, given that it is currently in chamber A, B, C or D, respectively.\par
\textbf{(i)} Explain why pA =\par
1 +\par
1 pB + 1\par
pD.\par
\textbf{(ii)} Determine pA .\par
\textbf{(iii)} Find the probability that a rabbit which starts in chamber A does not visit chamber C,\par
given that it eventually leaves the burrow through the tunnel in chamber A.\par
\end{flushleft}
\subsection*{Page 6}
\begin{flushleft}
STEP III, 2003 7\par
14 Write down the probability generating function for the score on a standard, fair six-faced die\par
whose faces are labelled 1, 2, 3, 4, 5, 6. Hence show that the probability generating function\par
for the sum of the scores on two standard, fair six-faced dice, rolled independently, can be\par
written as\par
\[1 t^2 (1 + t) 2 1 t + t^2 2 1 + t + t^2 2 .\]
36 -\par
Write down, in factorised form, the probability generating functions for the scores on two fair\par
six-faced dice whose faces are labelled with the numbers 1, 2, 2, 3, 3, 4 and 1, 3, 4, 5, 6, 8,\par
and hence show that when these dice are rolled independently, the probability of any given\par
sum of the scores is the same as for the two standard fair six-faced dice.\par
Standard, fair four-faced dice are tetrahedra whose faces are labelled 1, 2, 3, 4, the score being\par
taken from the face which is not visible after throwing, and each score being equally likely.\par
Find all the ways in which two fair four-faced dice can have their faces labelled with positive\par
integers if the probability of any given sum of the scores is to be the same as for the two\par
standard fair four-faced dice.\par
\end{flushleft}
\end{document}
